% !Mode:: "TeX:UTF-8"
% !TEX program = xelatex
%
% SlimeVR 神经网络残差校正器：当前训练模型、融合架构与工程状态
% 本文以仓库中的代码与 Phase 11/12 结果为准；研究假设与待实机验证项单独标注。

\documentclass[UTF8,11pt,a4paper]{ctexart}

\usepackage[margin=1in]{geometry}
\usepackage{amsmath,amssymb}
\usepackage{booktabs}
\usepackage{array}
\usepackage{graphicx}
\usepackage{float}
\usepackage{xurl}
\usepackage[hidelinks]{hyperref}
\usepackage{tikz}
\usetikzlibrary{positioning,arrows.meta}

\title{\LARGE\bfseries SlimeVR 误差校正模型\\[0.3em]\large 训练规范、双分支融合架构与当前实现状态}
\author{SlimeVR 研发记录}
\date{2026 年 10 月}

\begin{document}
\maketitle

\begin{abstract}
本文件记录当前 SlimeVR 下肢姿态校正研究的实际架构和工程边界。研究方案由两个职责分离的分支构成：Phase 6（P6）多流因果 CSTN 从四个下肢 IMU 与 HMD/VD 上身参考中预测短时姿态残差；Phase 9（P9）为四颗下肢 tracker 分别维护慢变世界偏航误差状态。Phase 11 的 Swing--Twist 合成器先应用 P9 的世界偏航修正，再只保留 P6 残差的非偏航 Swing 分量，以减少分支串扰。
截至 Phase 12，P6 已导出 ONNX 并接入 SlimeVR 服务端影子推理；P9 估计器和 P6/P9 合成仍是 Python 原型与离线验证，尚未集成到 Kotlin 运行链路。服务端默认影子模式，当前没有在线姿态修正。现有数值来自合成数据，尚不能证明实机长期漂移收益、真实 VD 条件下的融合增益或足底滑移改善。
\end{abstract}

\tableofcontents
\newpage

\section{范围、术语与状态口径}

\subsection{文档范围}
本文描述当前模型契约、网络结构、慢偏航估计器、正交融合公式、训练数据与验证状态。过往方案中提出但没有对应实现或验收证据的硬件滤波、视觉绝对锚点、ZUPT 锁脚、6 个下肢输出及 135 维输入，均不属于当前模型的已实现能力。

“P6”指 Phase 6 起源的短时神经残差分支；当前候选权重经 Phase 11 重训并施加 Yaw-Neutral 包装。“P9”指每 tracker 独立的偏航状态估计器。“P6+P9 融合”指 Phase 11 的离线 Swing--Twist 合成规则，不表示该融合已在服务端运行。

\subsection{当前状态摘要}
\begin{table}[H]
\centering
\small
\begin{tabular}{>{\raggedright\arraybackslash}p{2.4cm}>{\raggedright\arraybackslash}p{4.8cm}>{\raggedright\arraybackslash}p{6.6cm}}
\toprule
模块 & 当前实现 & 证据范围与限制\\
\midrule
P6 CSTN & PyTorch 训练、Yaw-Neutral 包装及 ONNX 单帧流式推理 & 合成动作扰动上有 Swing 残差改善；真实传感器失真匹配关系未知。\\
P9 状态估计 & Python 原型：四路 ZARU、步态候选观测与状态滤波 & Phase 12 的 A1 不接受原 PASS；模拟输入与服务端角速度来源不一致。\\
P6/P9 合成 & Python Swing--Twist 算子与单元测试 & 数学正交性通过；尚未在 Kotlin 服务端接入，融合收益没有独立确认。\\
SlimeVR 服务端 & P6 ONNX shadow bridge，60 Hz 采样 & 默认只记录候选值，不改 Avatar。P9 与融合器未接线；生产修正路径尚未完成。\\
硬件验收 & 未执行 & 当前没有目标实机与同步外部真值。\\
\bottomrule
\end{tabular}
\end{table}

\begin{figure}[H]
\centering
\begin{tikzpicture}[>=Stealth, node distance=0.8cm]
  \node[draw, rounded corners, text width=12.4cm, align=center, minimum height=0.85cm] (input) at (0,2.1)
    {124 维逐帧输入：HMD / 胸部 / 双上臂 + 双大腿 / 双小腿};
  \node[draw, rounded corners, text width=5.9cm, align=center, minimum height=1.1cm] (p6) at (-3.3,0)
    {P6 因果 CSTN\\四路短时 Swing 候选 + 门控 + 接触概率};
  \node[draw, rounded corners, text width=5.9cm, align=center, minimum height=1.1cm] (p9) at (3.3,0)
    {P9 独立状态估计\\四路偏航误差 \(e_i\)、变化率 \(b_i\) 与协方差};
  \node[draw, rounded corners, text width=12.4cm, align=center, minimum height=0.95cm] (fusion) at (0,-2.0)
    {Phase 11 离线合成：P9 世界 Y 偏航修正 + P6 Swing 投影/门控};
  \node[draw, rounded corners, text width=12.4cm, align=center, minimum height=0.85cm] (output) at (0,-3.55)
    {目标输出：下肢修正姿态送入 SlimeVR 骨架求解器};
  \draw[->, thick] (input.south west) -- (p6.north);
  \draw[->, thick] (input.south east) -- (p9.north);
  \draw[->, thick] (p6.south east) -- (fusion.north west);
  \draw[->, thick] (p9.south west) -- (fusion.north east);
  \draw[->, thick] (fusion) -- (output);
\end{tikzpicture}
\caption{研究态 P6/P9 融合结构。当前服务端只运行 P6 的影子推理，P9 和合成器尚未接入。}
\end{figure}

\section{输入契约与数据流}

\subsection{124 维输入 ABI}
模型每帧读取 8 路流，总计 124 个 float。上身四路各 17 维，下肢四路各 14 维。服务端在同一采样 tick 读取当前 tracker 值，但目前没有依据设备时间戳完成多源对齐。

\begin{table}[H]
\centering
\small
\begin{tabular}{c>{\raggedright\arraybackslash}p{3.1cm}c>{\raggedright\arraybackslash}p{8.2cm}}
\toprule
流索引 & 设备/身体通道 & 维度 & 字段顺序\\
\midrule
0 & HMD & 17 & 位置 3（服务端当前置零）、旋转 6D 6、线速度 3、角速度 3、有效 1、年龄 1。\\
1 & 胸部参考 & 17 & 与 HMD 相对的位置 3、旋转 6D 6、线速度 3、角速度 3、有效 1、年龄 1。\\
2--3 & 左/右上臂参考 & 17 each & 与 HMD 相对的位置 3、旋转 6D 6、线速度 3、角速度 3、有效 1、年龄 1。\\
4--7 & 左/右大腿、左/右小腿 & 14 each & 旋转 6D 6、加速度 3、角速度 3、有效 1、年龄 1。\\
\bottomrule
\end{tabular}
\caption{当前 v2 流式输入布局：\(4\times17+4\times14=124\)。}
\end{table}

6D 旋转采用旋转矩阵前两列，避免欧拉角不连续；四元数 ABI 使用 \([w,x,y,z]\)。无效或缺失的胸部、手臂及腿部流以零值和无效标志表示。当前服务端对 HMD 要求有效；胸部和手臂位置使用相对 HMD 的坐标。年龄字段目前写为零，不能视为已完成时延估计。

模型训练的腿部加速度来自合成数据中的线加速度。服务端在同一字段写入 \texttt{Tracker.getAcceleration()}。训练角速度由仿真直接生成；运行时优先读取 tracker 角速度字段，字段不可用时再由连续姿态四元数差分推导。Phase 12 审计未发现实际生产代码写入原始 tracker 角速度的调用点，因此运行时角速度与合成训练输入存在链路差异。

\subsection{服务端采样位置}
服务端在所有 \texttt{Tracker.tick()} 之后、\texttt{HumanPoseManager.update()} 之前采样。桥接器刷新 tracker 绑定、拼装 124 维帧、执行 ONNX 推理并缓存候选残差、门控和触地概率。Reset 和 tracker 断连会清空因果状态。输入字段及 Kotlin 映射见 \texttt{CorrectorBridge.kt}、\texttt{CorrectorConfig.kt} 与 \texttt{OnnxCorrectorEngine.kt}。

\section{P6：短时多流因果 CSTN}

\subsection{网络结构}
每个输入流独立嵌入到 96 维：线性层、LayerNorm、LeakyReLU。8 路嵌入分别进入一个单层因果 LSTM（hidden size 96），每路维护独立的 \(h,c\) 状态。每帧将 8 个 LSTM token 送入 3 层、8 头的 Transformer Encoder，进行 tracker 之间的空间交互。Transformer 之后的 token 均值形成全局上下文；它只用于构造上下文特征，不会替代腿部专属 token。

每个腿部姿态头将该腿 token 与全局上下文拼接，经独立 MLP 输出 3 维旋转向量，再转成单位四元数残差。四个腿部各有独立姿态头和门控头。残差向量幅度受最大 8 度限制。触地头从全局上下文输出左右脚两个概率。具体结构见 \texttt{ai\_corrector/corrector/model.py}。

因果时间信息由逐流 LSTM 状态跨帧递推；Transformer 只在当前帧的 8 路 token 之间注意，不沿时间轴做全序列注意。运行时逐帧输入、连续传递状态。这提供时序记忆，但当前训练与验收没有证明它能可靠保留或估计 30 分钟尺度的偏航误差。

\subsection{Yaw-Neutral 姿态残差}
基础 CSTN 输出局部右乘残差 \(\delta q_{\mathrm{local}}\)。Yaw-Neutral 包装器使用输入中的原始姿态把该增量共轭到世界系，并绕世界竖直轴 \(Y\) 分解：
\[
\delta q_{\mathrm{world}}
=\delta q_{\mathrm{swing}}\otimes\delta q_{\mathrm{twist},Y}.
\]
输出只保留 Swing，并转换回局部右乘约定。世界 \(Y\) 轴 Twist 被丢弃，避免 P6 与慢偏航分支重复修正同一自由度。网络本身分别输出 Swing 残差和门控值；离线合成器在门控低于 0.3 时使用恒等修正，高于门限时按门控权重做 identity-to-correction SLERP。导出器输出四路局部 Swing 四元数、四路门控、双脚接触概率以及下一帧的递归状态。

\subsection{训练损失}
Phase 11 Yaw-Neutral CSTN 使用以下加权目标：
\[
\mathcal{L}
=10\mathcal{L}_{\mathrm{swing}}
+15\mathcal{L}_{\mathrm{twist}}
+\mathcal{L}_{\mathrm{gate}}
+\mathcal{L}_{\mathrm{contact}}.
\]
\(\mathcal{L}_{\mathrm{swing}}\) 监督非偏航姿态残差；\(\mathcal{L}_{\mathrm{twist}}\) 惩罚世界竖直轴扭转；门控与接触头使用二元交叉熵。序列 mask 忽略填充帧。损失和包装逻辑见 \texttt{yaw\_neutral\_cstn.py}。

\section{P9：每 tracker 独立的慢偏航状态估计}

\subsection{状态定义}
P9 为左/右大腿和左/右小腿分别维护：
\[
\mathbf{x}_i =
\begin{bmatrix} e_i \\ b_i \end{bmatrix},
\qquad \mathbf{P}_i\in\mathbb{R}^{2\times2}.
\]
其中 \(e_i\) 是世界竖直轴的偏航姿态误差，\(b_i\) 是该姿态误差的变化率；它不是直接读取的芯片原始陀螺零偏。预测阶段执行 \(e_i\leftarrow\operatorname{wrap}(e_i+b_i\Delta t)\)，并传播协方差；每路独立更新并限幅。原型对 reset 和较大时间跳变实现状态复位；掉线期间如何恢复状态仍需服务端级验证。

\subsection{候选观测器}
\begin{itemize}
  \item \textbf{静止段 ZARU：}逐 tracker 检查低角速度、短窗稳定性和姿态突变。达到持续静止时长后，从静止区间的姿态变化估计偏航变化率，向对应状态通道提供率观测和几何对齐信息。
  \item \textbf{直线步态锚点：}在 HMD/VD 参考有效、HMD 水平位移达到阈值且大腿具有周期性摆动时，比较大腿前向轴水平投影与 HMD 行进方向，产生低权重偏航角观测。转身、原地运动、非直线行走或角度差过大时拒绝更新。
  \item \textbf{状态滤波：}ZARU 主要约束偏航误差率 \(b_i\)，步态观测只更新偏航角 \(e_i\)，避免弱步态观测把角度噪声耦合到长期漂移率。
\end{itemize}
这些是 Python 原型的工作逻辑，不代表已经有经过实机标定的观测噪声或完整生产实现。代码见 \texttt{bias\_estimator\_v3.py} 和 \texttt{bias\_state\_v3.py}。人体可以长时间保持非标准静态姿势，单纯对姿态角做低通会把真实姿势当成漂移；P9 的设计因此尝试依赖静止段的姿态变化率和严格条件下的步态锚点，而不把静态姿态本身拉回某个“标准朝向”。

\section{Phase 11：P6/P9 正交融合算子}

对 tracker \(i\)，令输入原始姿态为 \(R_{\mathrm{raw},i}\)，P9 估计的偏航误差为 \(\hat e_i\)，P6 输出局部残差为 \(\delta R_{6,i}\)，门控为 \(g_i\)。离线合成流程如下：
\begin{align}
R_{9,i} &= R_y(-\hat e_i)R_{\mathrm{raw},i},\\
R_{6,\mathrm{cand},i} &= R_{9,i}\delta R_{6,i},\\
D_i &= R_{6,\mathrm{cand},i}R_{9,i}^{T}.
\end{align}
将 \(D_i\) 转成四元数并沿世界 \(Y\) 轴作 Swing--Twist 分解。令其 Swing 分量为 \(q_{\mathrm{swing},i}\)，丢弃 Twist 分量。若 \(g_i<0.3\)，令门控 Swing 为单位旋转；否则令
\[
q_{\mathrm{swing},i}^{g}
=\operatorname{SLERP}(1,q_{\mathrm{swing},i},g_i).
\]
最终姿态和服务端局部右乘等价增量为
\[
R_{\mathrm{out},i}=R(q_{\mathrm{swing},i}^{g})R_{9,i},
\qquad
\delta R_{\mathrm{server},i}=R_{\mathrm{raw},i}^{T}R_{\mathrm{out},i}.
\]
因此慢偏航修正来自 P9；P6 只能增加通过门控的 Swing 修正，其世界 \(Y\) Twist 不参与合成。合成器的实现与数值边界测试见 \texttt{hybrid\_composer.py} 和 \texttt{test\_h1\_orthogonal\_operator.py}。

\section{训练数据与适用边界}

\subsection{P6 数据}
训练以 AMASS 动作序列构建合成 tracker 流。Phase 11 数据生成器记录 796 条训练序列、150 条验证序列和 190 条 FinalConfirm 序列。下肢各路独立注入动作相关 Swing 扰动；Phase 11 使用的扰动轴和幅度由角速度启发式构造（包含 \(0.040\|\boldsymbol{\omega}\times\mathbf{u}\|\) 项），它不是从真实 SlimeVR AHRS 误差标定或重放得到。VD/HMD 参考在主要实验中接近理想 AMASS 观测，不能覆盖真实噪声、延迟、安装偏差与重定位。

数据按序列隔离，不代表受试者级隔离。Phase 12 对 A2 复用 Phase 11 的 190 条 FinalConfirm，因此 A2 数字只作为探索性结果，不能作为新的确认集通过。

\subsection{评测结果及解释}
\begin{itemize}
  \item Phase 11 合成 30 分钟会话中，P9 ZARU-only 的平均端点偏航净收益为 \(+141.7188^\circ\)，候选 P9+P6 为 \(+141.7191^\circ\)。两者几乎相同，未显示 P6 的额外长期偏航收益；结果来自 10 条合成会话，Phase 12 已判定 A1 门禁无效且观测链路不一致，不能视为实机收益。
  \item Phase 11 的 Swing 确认性分析以真值构造的 P9 oracle Twist 作基线，平均 Swing 残差由 \(2.1799^\circ\) 降至 \(0.7261^\circ\)，报告改善 \(66.7\%\)。扰动为合成启发式、上身参考理想，且该 190 序列后来在 Phase 12 A2 被复用。因此它只说明模型能改善该合成目标，不能代表在线 P9 条件或真实 VD/IMU 表现。
  \item Phase 12 A2 报告显示 P9-only 与 Raw 相同，P9+P6 与 P6-only 相同。旧 FinalConfirm 被复用、简化 FK/脚滑代理和未纳入最终门禁的压力条件，使 A2 被裁定为探索结果，未证明融合增益。
\end{itemize}

\section{ONNX 与 SlimeVR 服务端接入状态}

\subsection{模型 ABI 与时延}
ONNX Runtime 使用单帧张量 \(x\in\mathbb{R}^{1\times1\times124}\)，输入递归状态 \(h,c\in\mathbb{R}^{8\times1\times1\times96}\)，输出四路残差 \(dq\)、四路门控 \(\mathrm{gate}\)、双脚接触概率 \(\mathrm{contact}\) 以及下一帧状态 \(h_n,c_n\)。模型以 60 Hz 逐帧运行，CPU 单线程 ONNX Session 常驻。Phase 12 本机测试在 108,000 逻辑帧上逐帧推理；3,000 个组装加推理样本的 P95 为 0.6571 ms。该计时未覆盖完整 \texttt{VRServer.tick()}，也不等于完整端到端延迟。

\subsection{当前服务端行为}
\begin{enumerate}
  \item \texttt{CorrectorConfig} 默认为 \texttt{enabled=false}、\texttt{shadowMode=true}。影子模式运行推理并缓存候选输出，不修改骨架姿态。
  \item Kotlin 接入只实现 P6 ONNX 推理；没有调用 Python P9，也没有把 Phase 11 合成算子移植进服务端。
  \item 当前 \texttt{applyCorrectionToTrackers()} 只计算局部变量 \texttt{corrected}，没有将它写回 tracker 或骨架求解器。因此实际在线姿态修正路径尚未完成；不可将启用配置等同于完整模型已生效。
  \item ONNX 引擎实现了 \texttt{close()}，但 Phase 12 复核未发现 \texttt{VRServer} 关闭生命周期调用桥接器关闭方法；Session 资源释放仍待补齐。
\end{enumerate}
完整本机验证状态见 \texttt{public/phase-12-prehardware-validation/results.md}。13 项指定 Kotlin 测试通过，S2 回放逐帧推理和延迟结果通过其局部门槛；Phase 12 整体验收尚未通过。

\section{安全设计、验证缺口与下一步}

网络门控只表示训练分布中的候选残差可信度，不是传感器真实度量或安全证明。上线前需要将低门控处理、异常输入旁路、角度/速率限幅、Reset 状态同步和可观测日志放在可审查的运行时路径中。短时 Swing 修正与长时偏航修正应分别报告，不能把平均姿态误差当成 30 分钟端点收益。

\begin{enumerate}
  \item \textbf{修正观测链路：}明确 SlimeVR 服务端可用的角速度、加速度、时间戳和质量字段；使训练与运行时特征的单位、坐标系、因果差分方式一致。
  \item \textbf{重新验证 P9：}修复静止姿势保持的独立门禁，验证静态姿势不会被冲洗；以新隔离序列分别测量四路漂移方向、端点误差、最差会话和置信区间。
  \item \textbf{确认融合增量：}使用未被训练、选模或 Phase 11/12 打开的新 source/subject 隔离集，将真实 P9 估计状态作为 P6 输入条件；报告 P9-only、P6-only、P9+P6 的配对端点和短时姿态结果。
  \item \textbf{补齐服务集成：}把 P9 状态机与正交合成器实现为服务端可测模块；将最终姿态安全写入骨架求解输入；补上 Session 生命周期关闭、Tracker 映射、时间同步和完整 \texttt{VRServer.tick()} 回放。
  \item \textbf{实机验证：}硬件可用后采集 IMU/VD/服务端输出，并同步外部光学或其他独立真值。单靠 AMASS 合成数据无法确认真实 AHRS 漂移模型或实际足部滑移改善。
\end{enumerate}

\section{实现文件索引}
\begin{itemize}
  \item P6 主干：\texttt{ai\_corrector/corrector/model.py}
  \item Yaw-Neutral 训练与包装：\texttt{ai\_corrector/corrector/yaw\_neutral\_cstn.py}
  \item Phase 9 状态估计：\texttt{ai\_corrector/corrector/bias\_estimator\_v3.py}、\texttt{bias\_state\_v3.py}
  \item Phase 11 融合算子：\texttt{ai\_corrector/corrector/hybrid\_composer.py}
  \item 服务端桥接与 ONNX：\texttt{server/server/core/src/main/java/dev/slimevr/tracking/processor/corrector/}
  \item 当前全局阶段索引：\texttt{public/README.md}
\end{itemize}

\end{document}
